
\appendix

% =====================================================================
\section{The matrices displayed}
\label{app:matrices}
% =====================================================================

This appendix displays the matrices of Sec.~\ref{sec:matrices} entry by entry, so that a reader can
check any statement about them by hand or with a few lines of code. In the displays a plus sign
stands for $+1$, a minus sign for $-1$ and a dot for $0$; the lines separate the $8\times8$ blocks.
The same matrices are available in machine-readable form \cite{SpacetimesRepo}.

%%<1671>
\emph{The $8\times8$ matrices $\tau_{1644}^A$ used in Eq.~(\ref{eq:tau}).} Let $\epsilon_{hpq4}$ be
the sign of the permutation $(h,p,q,4)$ of $(1,2,3,4)$, and zero if two of the indices coincide. For
$h=1,2,3$ and $p,q=1,\dots,4$ define the antisymmetric $4\times4$ matrices
\begin{align}
(s_h)_{pq}&=\epsilon_{hpq4}-\big(\delta_{p4}\delta_{qh}-\delta_{ph}\delta_{q4}\big),\nonumber\\
(t_h)_{pq}&=\epsilon_{hpq4}+\big(\delta_{p4}\delta_{qh}-\delta_{ph}\delta_{q4}\big),
\label{eq:st}
\end{align}
and, in $4\times4$ blocks,
\begin{align}
\tau_{1644}^h&=\begin{pmatrix}0&s_h\\ s_h&0\end{pmatrix},\qquad
\tau_{1644}^{7-h}=\begin{pmatrix}0&t_h\\ -t_h&0\end{pmatrix}\quad(h=1,2,3),\nonumber\\
\tau_{1644}^7&=\tau_{1644}^1\tau_{1644}^2\tau_{1644}^3\tau_{1644}^4\tau_{1644}^5\tau_{1644}^6=\mathrm{diag}(-1_4,1_4).
\label{eq:tau44}
\end{align}
These seven matrices satisfy $\tau_{1644}^1\tau_{1644}^2\tau_{1644}^3=\sigma_{88}$ (check
\texttt{44-B1}). Equation (\ref{eq:tau}) builds the seven octonionic matrices $\tau_{1671}^j$ from
them, and Eqs.~(\ref{eq:gamma4}) and (\ref{eq:gamma}) the eight Dirac matrices of
Table~\ref{tab:gammas}.
%%</1671>

\begin{table*}
\caption{\label{tab:gammas}The real matrices of the \SIG{$4+4$}{$7+1$} spacetime. First and second
row: the eight Dirac matrices $\gamma_{\tg}^A$; each is a signed permutation matrix. \SIG{The first
four, which belong to the space-like coordinates $x0,\dots,x3$, are symmetric and square to $+1$; the
last four, which belong to the time-like coordinates $x4,\dots,x7$, are antisymmetric and square to
$-1$.}{The matrix $\gamma^4$ of the time-like coordinate is symmetric and squares to $+1$; the other
seven, which belong to the space-like coordinates, are antisymmetric and square to $-1$.} Third row:
the product $\omega_{\tg}$ of the eight Dirac matrices, and the matrices $\Cm_{\tg}$ and $\Cp_{\tg}$
of Sec.~\ref{sec:other}. \SIG{$\omega_{1644}$ is the chirality matrix; it is diagonal, so that the
chiral projectors $P^{L,R}=\tfrac12(1\mp\omega_{1644})$ select the upper and the lower eight
components.}{$\omega_{1671}$ squares to $-1$; the chirality matrix is $\chi_{1671}=i\,\omega_{1671}$,
and the chiral projectors $P^{L,R}=\tfrac12(1\mp\chi_{1671})$ are complex.}}
\setlength{\tabcolsep}{1pt}
\begin{ruledtabular}
\begin{tabular}{cccc}
$\gamma_{\tg}^0$ & $\gamma_{\tg}^1$ & $\gamma_{\tg}^2$ & $\gamma_{\tg}^3$ \\[2pt]
{\fontsize{6}{7.2}\selectfont$\displaystyle{{MATRIX:gamma@SIG@[0]}}$} &
{\fontsize{6}{7.2}\selectfont$\displaystyle{{MATRIX:gamma@SIG@[1]}}$} &
{\fontsize{6}{7.2}\selectfont$\displaystyle{{MATRIX:gamma@SIG@[2]}}$} &
{\fontsize{6}{7.2}\selectfont$\displaystyle{{MATRIX:gamma@SIG@[3]}}$} \\[14pt]
$\gamma_{\tg}^4$ & $\gamma_{\tg}^5$ & $\gamma_{\tg}^6$ & $\gamma_{\tg}^7$ \\[2pt]
{\fontsize{6}{7.2}\selectfont$\displaystyle{{MATRIX:gamma@SIG@[4]}}$} &
{\fontsize{6}{7.2}\selectfont$\displaystyle{{MATRIX:gamma@SIG@[5]}}$} &
{\fontsize{6}{7.2}\selectfont$\displaystyle{{MATRIX:gamma@SIG@[6]}}$} &
{\fontsize{6}{7.2}\selectfont$\displaystyle{{MATRIX:gamma@SIG@[7]}}$} \\[14pt]
$\omega_{\tg}$ & $\Cm_{\tg}$ & $\Cp_{\tg}$ & \\[2pt]
{\fontsize{6}{7.2}\selectfont$\displaystyle{{MATRIX:omega@SIG@}}$} &
{\fontsize{6}{7.2}\selectfont$\displaystyle{{MATRIX:Cm@SIG@}}$} &
{\fontsize{6}{7.2}\selectfont$\displaystyle{{MATRIX:Cp@SIG@}}$} & \\
\end{tabular}
\end{ruledtabular}
\end{table*}

Table~\ref{tab:gammas} shows the eight Dirac matrices and three further matrices. In $8\times8$
blocks the chiral projectors of Eq.~(\ref{eq:PLPR}) are
%%<1644>
\begin{equation}
P_{1644}^{L}=\begin{pmatrix}1_8&0\\0&0\end{pmatrix},\qquad
P_{1644}^{R}=\begin{pmatrix}0&0\\0&1_8\end{pmatrix}.
\end{equation}
The Hermitian matrices of the quantized theory, $\beta_{1644}=-i\gamma_{1644}^4$ and
$\alpha_{1644}^j=-\gamma_{1644}^4\gamma_{1644}^j$ of Eq.~(\ref{eq:hamiltonian}), and the grading
$W_{1644}=i\,\Cm_{1644}\gamma_{1644}^4$ follow from the matrices displayed by one multiplication.
%%</1644>
%%<1671>
\begin{equation}
P_{1671}^{L}=\frac12\begin{pmatrix}1_8&i\,1_8\\-i\,1_8&1_8\end{pmatrix},\qquad
P_{1671}^{R}=\frac12\begin{pmatrix}1_8&-i\,1_8\\i\,1_8&1_8\end{pmatrix}.
\end{equation}
The real symmetric matrices of the Hamiltonian, $\beta_{1671}=\gamma_{1671}^4$ and
$\alpha_{1671}^j=\gamma_{1671}^4\gamma_{1671}^j$ of Eq.~(\ref{eq:alphabeta}), follow from the matrices
displayed by one multiplication.
%%</1671>

% =====================================================================
\section{The method of Cornwall, Jackiw and Tomboulis on a solvable model}
\label{app:toy}
% =====================================================================

The statements of Sec.~\ref{sec:method} about truncations of $\Gamma_2$ can be tested on a model in
which nothing is approximated: a ``field theory'' of a single real variable $\varphi$ with the action
\begin{equation}
S(\varphi)=\frac12\varphi^2+\frac{\lambda}{24}\varphi^4 .
\end{equation}
Its partition function $Z=\int d\varphi\,e^{-S}/\sqrt{2\pi}$ and its ``propagator''
$G=\langle\varphi^2\rangle$ are ordinary integrals. The CJT action (\ref{eq:CJT}) is a function of
one number,
\begin{align}
\Gamma(G)&=\frac12\ln\frac1G+\frac12(G-1)+\Gamma_2(G),\nonumber\\
\Gamma_2&=\frac{\lambda G^2}{8}-\frac{\lambda^2G^4}{48}+\frac{\lambda^3G^6}{48}-\frac{5\lambda^4G^8}{128}+\dots,
\label{eq:toy}
\end{align}
whose terms are the two-particle-irreducible diagrams with two, three, four and five loops. The gap
equation (\ref{eq:SD}) is $1/G=1+\Sigma$ with
$\Sigma=\lambda G/2-\lambda^2G^3/6+\lambda^3G^5/4-5\lambda^4G^7/8+\dots$, and the value of $\Gamma$
at its solution should equal $-\ln Z$. The coefficients follow from the Gaussian moments
$\langle\varphi^{2k}\rangle_0=(2k-1)!!$: the perturbation series of the exact propagator is
$G=1-\lambda/2+2\lambda^2/3-11\lambda^3/8+34\lambda^4/9-\dots$, and inverting it and expressing
$1/G-1$ through $G$ itself gives the series for $\Sigma$. A unit test of the program verifies the
four coefficients against the exact integral.

\begin{figure*}
\includegraphics[width=\textwidth]{fig_cjt_toy}
\caption{\label{fig:toy}The CJT method in zero dimensions. (a) The exact propagator, the solutions of
the gap equation with $\Gamma_2$ truncated at two, three and four loops, and ordinary perturbation
theory. (b) Difference between the effective action at the stationary point and the exact $-\ln Z$.}
\end{figure*}

\begin{table}
\caption{\label{tab:toy}The propagator $G=\langle\varphi^2\rangle$ of the zero-dimensional model: exact
value, and solutions of the gap equation with $\Gamma_2$ truncated at two to five loops. A dash means
that the truncated gap equation has no acceptable solution. The values of perturbation theory are
shown in Fig.~\ref{fig:toy}.}
\begin{ruledtabular}
\begin{tabular}{cccccc}
$\lambda$ & exact & 2 loops & 3 loops & 4 loops & 5 loops \\
\hline
{{ROWS:cjt_toy_table|cols=0-5}}
\end{tabular}
\end{ruledtabular}
\end{table}

Figure~\ref{fig:toy} and Table~\ref{tab:toy} teach three lessons, which carry over to field theory.
\emph{The two-loop truncation is uniformly good.} It stays within {{GEN:cjt_toy_hartree_err}} per
cent of the exact propagator for all couplings from 0.01 to 30, where the perturbation series is
useless beyond $\lambda\approx0.5$; its solution, $G=(\sqrt{1+2\lambda}-1)/\lambda$, has the correct
qualitative behavior for $\lambda\to\infty$. \emph{More loops help only at weak coupling.} The
series (\ref{eq:toy}) is asymptotic, like the perturbation series. At $\lambda=1$ the error of the
effective action falls from {{GEN:cjt_toy_gamma_err_2}} with two loops to
{{GEN:cjt_toy_gamma_err_4}} with four. The truncations at three and at five loops, whose last term
in $\Sigma$ is negative, have no acceptable solution for $\lambda\ge5$ and for $\lambda\ge2$,
respectively; the entry 1.0000 at $\lambda=3$ is a solution, but a poor one. \emph{The stationary
value is more accurate than the propagator,} because an error of first order in the propagator
changes the value of a stationary functional only in second order.

% =====================================================================
\section{Machine checks, tests and reproduction}
\label{app:checks}
% =====================================================================

\begin{table}
\caption{\label{tab:checks}The machine checks by group, with the number of checks for each spacetime.}
\begin{ruledtabular}
\begin{tabular}{c>{\raggedright\arraybackslash}p{0.66\columnwidth}cc}
group & what the checks verify & $4+4$ & $7+1$ \\
\hline
{{ROWS:checks_table}}
\end{tabular}
\end{ruledtabular}
\end{table}

\emph{Machine checks.} The statements of Secs.~\ref{sec:matrices} to \ref{sec:pairs} about matrices
are verified by {{CHECKS}} checks, {{GEN:n_checks_@SIG@}} of them for the spacetime of this paper.
An identifier has the form \texttt{44-C1} or \texttt{71-C1}: the signature, a letter for the group
of Table~\ref{tab:checks}, and a number. {{GEN:n_checks_exact}} checks are exact: they are
identities between integers and are proofs. The other {{GEN:n_checks_numerical}} involve
exponentials of matrices or eigenvalues and are verified in floating-point arithmetic, to tolerances
stated in the program. All checks pass. The log of the program lists every check with its complete
statement \cite{SpacetimesRepo}.


\emph{An independent second program.} A program can be wrong in the same way twice. A second
program, in another language and sharing no code with the first, therefore rebuilds the matrices
from their definitions and repeats the essential checks by different methods: linear independence
is proved with exact Gram matrices instead of ranks modulo a prime; the question which matrices
anticommute with all eight Dirac matrices is settled by enumerating the 256 products instead of
solving a linear system; and the classification of the pairs is redone in exact rational
arithmetic. It also compares the matrices, entry by entry, with those of the first program and with
the matrices printed in the notebooks from which the definitions were taken.

\emph{Tests.} The {{GEN:n_unit_tests}} unit tests of the main program compare each numerical
routine with a closed result or with an independent method: the exchange amplitudes with the closed
forms (\ref{eq:FSFP}), the gap equation with the textbook limits, the gradient of the CJT potential
with numerical derivatives, the two-particle-irreducible coefficients of Appendix~\ref{app:toy} with
the exact integral, the equations of state of the dark fluids with the conservation of energy, and
the Einstein-Boltzmann solver with stored reference values.

\emph{Reproduction.} From the root of the repository, the commands
\begin{verbatim}
cd code/spacetimes-rs
cargo build --release
cargo test --release
cd ../..
B=code/spacetimes-rs/target/release/spacetimes
for c in matrices channels gap dimer \
    symmetry cjt background spectrum \
    forecast superfluid halos
do $B $c; done
python3 code/scripts/\
verify_matrices_independent.py
python3 code/scripts/\
validate_cmb_against_class.py
python3 code/scripts/make_tables.py
python3 code/scripts/make_figures.py
python3 code/scripts/build_manuscripts.py
\end{verbatim}
compile and test the program, recompute every data file, repeat the independent verification and
the comparison with \textsc{class}, and rebuild the tables, the figures and both manuscripts. The complete run takes about an hour and a
half on a computer with two processor cores, most of it for the fits to the mock data.

% =====================================================================
\section{Numerical methods}
\label{app:num}
% =====================================================================

\begin{table}
\caption{\label{tab:cmbval}Differences, in per cent, between the spectra of this work, lensed with the
lensing potential of the reference, and those of \textsc{class}: largest and root-mean-square values
of $D_\ell/D_\ell^{\rm ref}-1$ for TT and EE and of
$(D_\ell^{TE}-D_\ell^{TE,\rm ref})/\sqrt{D_\ell^{TT}D_\ell^{EE}}$ for TE, where
$D_\ell\equiv\ell(\ell+1)C_\ell/2\pi$ and $C_\ell$ is an angular power spectrum. Below $\ell=30$ the
polarization depends on the description of reionization; above $\ell=2000$ the lensed result depends
on the continuation of the spectra beyond $\ell=2500$.}
\begin{ruledtabular}
\begin{tabular}{lcccccc}
 & \multicolumn{2}{c}{TT} & \multicolumn{2}{c}{TE} & \multicolumn{2}{c}{EE} \\
multipoles & max & rms & max & rms & max & rms \\
\hline
{{ROWS:cmb_validation_table}}
\end{tabular}
\end{ruledtabular}
\end{table}

All calculations are performed by one program, written in the language Rust, with
{{GEN:n_unit_tests}} unit tests; the figures and tables are produced from its output files by
scripts. This appendix summarizes the methods; the source code is commented, and the repository
\cite{SpacetimesRepo} documents, for every number of this paper, the command that produces it.

\emph{Exact algebra.} Matrices with integer entries are multiplied in 64-bit integer arithmetic.
Ranks are computed by elimination modulo a prime; since a rank modulo a prime cannot exceed the rank
over the rational numbers, a full rank found in this way proves linear independence. Complex
matrices are stored as pairs of integer matrices.

\emph{Symmetry and pairs.} The commutant of a set of matrices is the null space of a linear system,
found by Gaussian elimination with pivoting; the quantum numbers of particles and pairs are
eigenvalues of Casimir operators, computed in floating-point arithmetic with the Jacobi method and,
in the second program, in exact rational arithmetic.

\emph{The paired state.} The gap equation (\ref{eq:gap}) is discretized on a radial momentum grid of
Gauss-Legendre panels that are refined geometrically towards the Fermi momentum, and solved by
iteration; the critical temperature is found by bisection on the largest eigenvalue of the linearized
kernel. The CJT potential (\ref{eq:Vpair}) is evaluated on the same grid; its stationary point is
found by Newton's method with the analytic Jacobian. The radial Schr\"odinger equation of the bound
pair is integrated with the fourth-order Runge-Kutta method on a grid that is uniform in
$\ln(1+\rho/\rho_0)$, which resolves the singularity of the potential at the origin.

\emph{Matter, temperature and vacuum.} The Fermi integrals of Sec.~\ref{sec:normal} are closed
expressions. The thermal integrals of Sec.~\ref{sec:thermal} are tabulated and interpolated with
cubic polynomials. The ladder equation (\ref{eq:ladder}), linearized, is an eigenvalue problem on a
logarithmic momentum grid, solved by power iteration.

\emph{Cosmology.} The background is integrated by Gauss-Legendre quadrature. Recombination follows
Saha equilibrium at early times and then the rate equations of the effective three-level atoms
\cite{SeagerSasselovScott1999,SeagerSasselovScott2000,WongMossScott2008}; reionization is a step
whose redshift is fitted to the optical depth. The perturbations obey 50 coupled equations per
wavenumber, integrated with an adaptive Runge-Kutta method of Dormand and Prince, with the
tight-coupling approximation of first order at early times. The spectra are obtained by integration
along the line of sight with tabulated spherical Bessel functions. Table~\ref{tab:cmbval} gives the
comparison with \textsc{class} by ranges of the multipole.


\emph{Mock likelihood.} For Gaussian fields observed on a fraction $f_{\rm sky}$ of the sky,
\begin{align}
\chi^2={}&\sum_\ell(2\ell+1)\,f_{\rm sky}\Big[\Tr\big(\hat C_\ell C_\ell^{-1}\big)+\ln\frac{\det C_\ell}{\det\hat C_\ell}-n\Big]\nonumber\\
&+\Big(\frac{\tau-\tau_0}{\sigma_\tau}\Big)^2,
\label{eq:likelihood}
\end{align}
where $C_\ell$ is here the $n\times n$ matrix of the model spectra plus noise, for temperature and
polarization, $\hat C_\ell$ the same matrix for the data, and the last term the prior on the optical
depth $\tau$ (central value $\tau_0$, standard error $\sigma_\tau$). A data set is a list of disjoint patches of the sky, each with its own $f_{\rm sky}$, noise
and range of multipoles, and their $\chi^2$ add.

